\documentclass[english]{lni}
% for German use:
%\documentclass{lni}

\IfFileExists{latin1.sty}{\usepackage{latin1}}{\usepackage{isolatin1}}  % comment out this line for compilation with e.g. overleaf

\usepackage{graphicx}
\usepackage{fancyhdr}
\usepackage{listings} %if lstlistings is used
\usepackage{changepage} %for changing topmargin on first page
\usepackage[figurename=Fig., tablename=Tab., small]{caption}[2008/04/01]
\renewcommand{\lstlistingname}{List.}    % Listingname is now List. 

\fancypagestyle{titlepage}{
\fancyhead[RO]{\small N. Damer,  M. Gomez-Barrero,  K. Raja,  C. Rathgeb,  A.  Sequeira,\linebreak  M. Todisco,  and A. Uhl (Eds.): BIOSIG 2023, \linebreak Lecture Notes in Informatics (LNI), Gesellschaft f\"ur Informatik, Bonn 2023} % do NOT modify these lines
\fancyfoot{}}

%Beginning of page count for this paper
\setcounter{page}{1}

%head line settings
\pagestyle{fancy}
\fancyhead{} % clears the settings
\fancyfoot{} % clears footer settings
\renewcommand{\headrulewidth}{0.4pt} %horizontal line below header
\setcounter{footnote}{0}

\author{Firstname1 Lastname1\footnote{Faculty1, Department1, Address1, Place1, emailaddress1@author1} \, Firstname2 Lastname2\footnote{Faculty2, Department2, Address2, Place2, emailaddress2@author2} \ use the same notation for additional authors}
\title{Title of your paper which can have up to two lines}
\begin{document}

\maketitle

\renewcommand{\refname}{References}
\setcounter{footnote}{2} %Change to the number of authors for a correct numbering of the foot notes
\thispagestyle{titlepage}
%header setting after the second page
\pagestyle{fancy}
\fancyhead{} % clears header settings
\fancyhead[RO]{\small Short title of your work \hspace{25pt}  \hspace{0.05cm}}
\fancyhead[LE]{\hspace{0.05cm}\small  \hspace{25pt} Firstname1 Lastname1 and Firstname2 Lastname2}
\fancyfoot{} % clears all footer settings
\renewcommand{\headrulewidth}{0.4pt} %line below header

\begin{abstract}
	The \LaTeX-Class \texttt{lni} uses the layout for articles in LNI. This document describes the usage of the \LaTeX-Class and gives some examples. This abstract should give a short overview about your work and can have between 70 and 150 words. The formatting will be done automatically within the abstract area.
\end{abstract}
\begin{keywords}
	LNI Guidelines, \LaTeX Template. The formatting will be done automatically within the keywords-Range.
\end{keywords}

\section{Usage}
The GI gives specification \texttt{http://www.gi-ev.de/LNI} on the formatting of documents inside the LNI. For \LaTeX-Documents this is done with the help of the document class \texttt{lni}. This document gives information on version 0.6. with bibliography styles \texttt{lnig.bst} (deutsch) version 2.0 and \texttt{lni.bst} (english) version 1.0.

The document class can be included within the document header.

\begin{verbatim}
\documentclass{lni}
\end{verbatim}

The class inherits from the \texttt{article}-class and needs the Times-Roman font to run properly. Include \texttt{times} and \texttt{verbatim}-styles. These should be shipped with your \LaTeX-installation, and if not they can be downloaded freely on the internet.

You will need to include the styles fancyhdr.sty and changepage.sty also, which can be found inside the example package or under www.ctan.org .

Possible document options are: \texttt{english}. For choosing a language, see section~\ref{sprache}.\linebreak
\clearpage


For the correct naming of the captions use the following commands:

\begin{verbatim}
\usepackage[figurename=Fig.,tablename=Tab.,small]{caption}
[2016/04/22]
\renewcommand{\lstlistingname}{List.}  
\end{verbatim}
In English texts use \texttt{figurename} : \texttt{Fig.}.

The numbering of your article will be done by the editor. With the help of the command:
\begin{verbatim}
\setcounter{page}{11}
\end{verbatim}

you can set the start number of the page counter if you need to change it on your own.

If \texttt{lstlisting} is used you have to include listings.sty

Please make use of the default macros of the class \texttt{article}. If you have to define macros on your own, you should put them into \texttt{.sty}-files to ease to combination of articles later on.

\subsection{Language}
\label{sprache}

Both German and English are supported by this documents class. German is the default. For completely English written articles you need to set the document option \texttt{english}.

To switch language within the text, use \verb|\selectlanguage| from the Babel-Package. This package will be loaded by the \texttt{lni}-class -- there is no need to load it again.

\texttt{lni} uses the word separation rules of the new German spelling. The shortcut is \texttt{ngerman} (but not \texttt{german}). If your \LaTeX\ does not support \texttt{ngerman}, you have to change your system (see \texttt{www.dante.de}).

To write the complete text in English set the document option:

\begin{verbatim}
\documentclass[english]{lni}
\end{verbatim}

Then \LaTeX uses English words for fix text components (e.g.  `Literaturverzeichnis' to `References').

Therefore you also have to delete the command:
\begin{verbatim}
\renewcommand{\refname}{Literaturverzeichnis}
\end{verbatim}

As said before, do not forget to change the abbreviation of figures to \texttt{Fig.}:
\begin{verbatim}
\usepackage[figurename=Fig.,tablename=Tab.,small]{caption}
[2016/04/22]  
\end{verbatim}

Additional language specific classes can be used manually. This document uses for example:

\begin{verbatim}
\usepackage{latin1}
\end{verbatim}

\subsection{Title, Authors, Abstract and Keywords}

The title is stored in \verb|\title|. Authors are defined within \verb|\author|. The title itself can be set by \verb|\maketitle|. Information on the authors are given inside footnotes(Faculty, Address, Email) \verb|\footnote| within \verb|\author|.

The following lines will set to document information for this document:

%\begin{verbatim}
%\author{Firstname1 Lastname1 \footnote{...} \ and ...}
%\title{The title}
%\begin{document}
%\maketitle

% Correct naming for German texts - outlined for English text
%\renewcommand{\refname}{Literaturverzeichnis}

%Reset the counter for footnotes because of multiple author footnotes
\setcounter{footnote}{2}
\begin{abstract}
	Abstract Text.
\end{abstract}
\begin{keywords}
	Keywords, Text.
\end{keywords}
\end{verbatim}

The command \verb|\renewcommand{\refname}{Literaturverzeichnis}| sets the title of the bibliography to "Literaturverzeichnis". For English documents the command has to be removed to get "References".

The command \verb|\setcounter{footnote}{2}| resets the counter of the footnotes for cases the multiple authors with footnotes are used.

\subsection{Text structure}

A variety of parameters for text structures were redefined within the \texttt{lni}-class. To apply these changes you will need to use the exact same commands.

The layout for flowing text is the defaule.
Section titles are set with \verb|\section|, subsections with \verb|\subsection|. Footnotes can be set with \verb|\footnote|.\footnote{Layout is defined within the document class}
\begin{itemize}
	\item Enumerations are done help of the \verb|\itemize| environment.
	\item The formatting is predefined within the \texttt{lni}-class.
	      \begin{itemize}
		      \item Enumerations should have at least 2 levels.
	      \end{itemize}
\end{itemize}
\begin{enumerate}
	\item Numbered enumerations are done help of the \verb|\enumerate| environment.
	      \begin{enumerate}
		      \item The formatting is predefined within the \texttt{lni}-class.
	      \end{enumerate}
\end{enumerate}

\subsection{Figures}

Figures are included with default \LaTeX-commands. You have to insert the caption \verb|\caption| right after the figure definition but before the \verb|\centering|-command to ensure correct centering of the figure. (Please avoid using \verb|\begin{center}...|\verb|\end{center}| in the context of figures which can lead to an unwanted output.

The positioning of the figure should be done by \texttt{[htb]}. Figure ~\ref{logo} is an example for that.

\begin{verbatim}
\begin{figure}[htb]
\centering
    \includegraphics[width=2cm]{gilogo}
    \caption{\label{logo}Icon of GI}
\end{figure}
\end{verbatim}

\begin{figure}[htb]
	\centering
	\includegraphics[width=2cm]{gilogo}
	\caption{\label{logo}Icon of GI}
\end{figure}

If possible, use the setting \texttt{[h]}, to allow \LaTeX to shift your picture automatically, instead of using \texttt{t-top} for the beginning or \texttt{b-bottom} for the bottom of the page to avoid unwanted spacing above and below the figure.

Other figure styles as \texttt{subfigure} have not been tested yet.

\subsection{Tables}
Tables are done with default \LaTeX-commands. Insert the table caption \verb|\caption| right after the table definition and use \verb|\centering| before the definition.
\begin{verbatim}
\begin{table}[htb]
\centering
\begin{tabular}{l l l}
\hline
Heading Layers & Example & font size and type\\
\hline
Title (left justified) & The title... & 14 pt, bold\\
\hline
Heading 1 & 1 Introduction & 12 pt, bold\\
\hline
Heading 2 & 2.1 Title & 10 pt, bold\\
\hline
\end{tabular}
\caption{\label{tab1}The types of headings}
\end{table}
\end{verbatim}

Aligning tables is done with \verb|[htb]|. Table~\ref{tab1} shows an example for that.

\begin{table}[htb]
	\centering
	\begin{tabular}{l l l}
		\hline
		Heading              & Example        & Font size and type \\
		\hline
		Title (left aligned) & The title...   & 14 pt, bold        \\
		Heading 1            & 1 Introduction & 12 pt, bold        \\
		Heading 2            & 2.1 Title      & 10 pt, bold        \\
		\hline
	\end{tabular}
	\caption{\label{tab1}The headings}
\end{table}

\subsection{Program code}
The LNI template demands you to indent the listing. Therefore the \texttt{lni} document class defines the \texttt{verbatim}-environment.
\begin{verbatim}
public class Hello { 
    public static void main (String[] args) { 
        System.out.println("Hello World!"); 
    } 
} 
\end{verbatim}

As an alternative you can also use the \texttt{lstlisting}-environment.
\begin{verbatim}
	\lstset{basicstyle=\ttfamily}	%defines the font family
	\begin{lstlisting} [captionpos=b, caption=Description, label=L1,
	 xleftmargin=0.5cm]
	 %Description and Lable can be defined here.
	...
	\end{lstlisting}
\end{verbatim}

Listing~\ref{L1} shows an example, the was created with the \texttt{lstlisting}-environment.

\lstset{basicstyle=\ttfamily}
\begin{lstlisting} [captionpos=b, caption=Description, label=L1, xleftmargin=0.5cm]
public class Hello { 
    public static void main (String[] args) { 
        System.out.println("Hello World!"); 
    } 
}
\end{lstlisting}

\subsection{Formulas and Equations}

The environments \texttt{equation} and \texttt{eqnarray} ensure correct indent and numbering of formulas.

\begin{verbatim}
	\begin{equation}
	  1=4-3
	\end{equation}

	\begin{eqnarray}
	  2=7-5\\
	  3=2-1
	\end{eqnarray}
\end{verbatim}

creates...
\begin{equation}
	1=4-3
\end{equation}
and
\begin{eqnarray}
	2=7-5\\
	3=2-1
\end{eqnarray}

\subsection{Bibliography}
The \texttt{lni}-class comes with an integrated \texttt{lnig}-bibliography style. There is no need to additionally add \verb|\bibliographystyle{lni}|. You can add the bibliography as follows:

\begin{verbatim}
\bibliography{lniguide}
\end{verbatim}

When using the document option \texttt{english} for English text, the \texttt{lni}-class automatically uses the English bibliography style. If you manually switch between languages within the document, the \texttt{lni}-class selects the style for the language active, when calling \verb|\bibliography|.

\subsubsection{References - Citation}

Formatting and abbreviations are done automatically for \texttt{book}, \texttt{inbook}, \texttt{proceedings}, \texttt{inproceedings}, \texttt{article}, \texttt{online} and \texttt{misc}. You will find examples in the \texttt{lniguide.bib} file inside the template folder.

Additional references or fields need to be edited manually.

\subsubsection{Examples}
The last chapter gives an example on a bibliography for:
\begin{itemize}
	\item a book with one author - \cite{Ez10}
	\item ... two authors - \cite{AB00}
	\item an article within a proceeding, written by three authors - \cite{ABC01}
	\item an article within LNI, written by more than authors - \cite{Az09}
	\item two books of the same year, written by the identical four authors - \cite{Wa14} and \cite{Wa14b}
	\item a journal - \cite{Gl06}
	\item a website - \cite{GI14}
	\item additional literature without concrete information on the authors \cite{XX14}
\end{itemize}

\section{Header and page numbering}
The first page has its own already defined header and differs from the odd an even headers. You must not modify the header on the first page.

It ensures the correct formatting of the headers from the second page on.

The headers of the even pages contain the page number and authors on the left. Up to three authors are supported here. For more then three authors use the first author's name and et al. .
Too long author lists can be shortened by the editor.

Headers of odd pages contain the page number and the (short) title on the right.

\section{Installation}

\texttt{lni.cls} and \texttt{lni.bst} need to be accessible by your \TeX-System. The class uses the packages \texttt{times} and \texttt{verbatim}, which are normally included in the \LaTeX-Installation or have to be installed otherwise (free source: \texttt{www.dante.de}).

\bibliographystyle{lnig}
\bibliography{lniguide}

\end{document}
